\documentclass[12pt,a4paper]{article}
\usepackage[utf8]{inputenc}
\usepackage[indonesian]{babel}
\usepackage{amsmath}
\usepackage{amsfonts}
\usepackage{amssymb}
\usepackage{geometry}

\geometry{margin=2.5cm}

\title{Menghitung Luas Permukaan Benda Putar Sumbu-x}
\author{Ahmad Fakhrul Bawani}
\date{}

\begin{document}

\maketitle

\section*{Soal}

Find the area of the surface generated by revolving the curve:
\begin{equation*}
  y = 6\sqrt{x}, \quad \text{untuk } 27 \le x \le 91, \quad \text{mengelilingi sumbu-}x
\end{equation*}

\subsection*{1. Rumus Umum}
Rumus luas permukaan ($S$) benda putar yang dihasilkan oleh perputaran mengelilingi sumbu-$x$ dari $x = a$ sampai $x = b$ didefinisikan sebagai:
\begin{equation*}
  S = 2\pi \int_{a}^{b} y \cdot \sqrt{1 + \left(\frac{dy}{dx}\right)^2} \, dx
\end{equation*}

\subsection*{2. Menentukan Turunan Fungsi}
Tulis ulang fungsi ke dalam bentuk pangkat: $y = 6x^{\frac{1}{2}}$. Cari turunan pertama $y$ terhadap $x$:
\begin{equation*}
  \frac{dy}{dx} = 6 \cdot \left(\frac{1}{2}x^{-\frac{1}{2}}\right) = 3x^{-\frac{1}{2}} = \frac{3}{\sqrt{x}}
\end{equation*}

Kuadratkan hasil turunan tersebut:
\begin{equation*}
  \left(\frac{dy}{dx}\right)^2 = \left(\frac{3}{\sqrt{x}}\right)^2 = \frac{9}{x}
\end{equation*}

Masukkan ke dalam komponen akar panjang busur:
\begin{equation*}
  \sqrt{1 + \left(\frac{dy}{dx}\right)^2} = \sqrt{1 + \frac{9}{x}}
\end{equation*}

\subsection*{3. Solusi Integrasi}
Substitusikan komponen $y = 6\sqrt{x}$ dan bentuk akar di atas ke dalam rumus luas permukaan dengan batas $a = 27$ dan $b = 91$:
\begin{align*}
  S &= 2\pi \int_{27}^{91} (6\sqrt{x}) \cdot \sqrt{1 + \frac{9}{x}} \, dx \\
  S &= 12\pi \int_{27}^{91} \sqrt{x \cdot \left(1 + \frac{9}{x}\right)} \, dx \\
  S &= 12\pi \int_{27}^{91} \sqrt{x + 9} \, dx \\
  S &= 12\pi \int_{27}^{91} (x + 9)^{\frac{1}{2}} \, dx
\end{align*}

Gunakan metode substitusi atau integrasi langsung:
\begin{align*}
  \text{Misalkan: } & u = x + 9 \implies du = dx \\
  \text{Batas baru: } & \text{Untuk } x = 27 \implies u = 27 + 9 = 36 \\
  & \text{Untuk } x = 91 \implies u = 91 + 9 = 100
\end{align*}

Substitusikan variabel baru $u$ ke dalam integral:
\begin{align*}
  S &= 12\pi \int_{36}^{100} u^{\frac{1}{2}} \, du \\
  S &= 12\pi \left[ \frac{2}{3}u^{\frac{3}{2}} \right_{36}^{100} \\
  S &= 8\pi \left[ u\sqrt{u} \right_{36}^{100} \\
  S &= 8\pi \left( \left[ 100\sqrt{100} \right] - \left[ 36\sqrt{36} \right] \right) \\
  S &= 8\pi \left( 100(10) - 36(6) \right) \\
  S &= 8\pi (1000 - 216) \\
  S &= 8\pi (784) \\
  S &= 6272\pi
\end{align*}

Jadi, luas permukaan benda putar tersebut adalah \textbf{$6272\pi$ satuan luas}.

\end{document}